\chapter{Arguments, exponentielle imaginaire et trigonométrie}
\section*{Exercice 26.1 -- Forme trigonométrique}
[Base] Écrire $-1+i\sqrt3$ sous forme trigonométrique et exponentielle.

\section*{Exercice 26.2 -- Puissance}
[Méthode] Calculer $(\sqrt3+i)^6$.

\section*{Exercice 26.3 -- Addition trigonométrique}
[Preuve] Déduire $\cos(a+b)$ et $\sin(a+b)$ de $e^{i(a+b)}=e^{i a}e^{i b}$.

\section*{Exercice 26.4 -- Linéarisation}
[Méthode] Exprimer $\cos^3x$ comme combinaison de cosinus d'angles multiples.

\section*{Exercice 26.5 -- Somme de phases}
[Approfondissement] Pour $e^{i t}\ne1$, calculer $\sum_{k=0}^{n-1}e^{i kt}$. Que se passe-t-il si $t$ est multiple de $2\pi$ ?

\section*{Exercice 26.6 -- Argument d’un produit}
[Critique] Les arguments principaux de deux nombres sont $3\pi/4$ et $3\pi/4$. L'argument principal du produit est-il $3\pi/2$ dans $]-\pi,\pi]$ ?
